% MAIN TABLE -- zero-shot text-to-video R@1 on five benchmarks.
%
% SINGLE CONFIGURATION. One SCA row. Query-weighted centroid, no frame slots, tau_w = 0.1,
% LoRA r8, batch 128, lr 2e-5, 5 epochs (arm T9). Evaluated on every benchmark and reported
% everywhere, including where it trails. The configuration was fixed before the transfer
% numbers were seen -- the simplest setting the ablation supports -- not chosen by which
% benchmark a variant happened to win.
%
% n = 3 SEEDS. The SCA row is the mean of three runs of that configuration differing only in
% run_cfg.seed (t9, s1=51, s2=52), with the sample standard deviation. Per-benchmark seed
% values, for the appendix:
%
%   MSR-VTT     54.8 / 54.7 / 54.4     DiDeMo      51.5 / 51.0 / 51.2
%   ActivityNet 55.8 / 56.1 / 55.9     VATEX       90.5 / 90.5 / 90.7
%   AudioCaps   35.2 / 34.5 / 35.2
%
% EVERY CELL VERIFIED. scripts/audit_eval_geometry.py compares each result directory's own
% recorded config against its training arm's, across all eval roots. It found 25 cells scored
% under a geometry they were never trained with -- t12 given frame slots it never learned, t8
% and t11 scored at tau_w 0.1 rather than the 0.05 and 0.2 they trained at -- and those were
% re-run. No cell behind this table was among them, and the audit now exits clean on
% workdir/e1_frames. That check runs in the harvest, so the table cannot silently drift back.
%
% Convention: $^{\S}$ published by the authors; $^{\ast}$ measured by us in one environment.
% GRAM released its weights, so the same checkpoint appears both ways -- and the difference
% between those two rows IS the environment shift ($-2.3$ MSR-VTT, $-3.5$ DiDeMo, $+6.5$
% VATEX). That is why $^{\S}$ and $^{\ast}$ rows are not comparable in absolute terms, and why
% the $^{\ast}$ GRAM row is the honest baseline for SCA.
%
% HyperGRAM is cited from its paper. Our reimplementation is not a row here: no code is
% released, the hyperbolic branch admits two readings, and the arm at the paper's own learning
% rate is still training. See experiments/results/HYPERGRAM_STATUS.md -- and do not quote the
% 37.4 from gram_hyp2, which was trained at lr 2e-5.
\begin{table*}[t]
\centering
\caption{Zero-shot text-to-video retrieval, R@1. $^{\S}$ as published; $^{\ast}$ measured by us
in a single environment. GRAM is the only method here that released weights, so the same
checkpoint appears as both rows; their difference is a \emph{measured} environment shift
($-2.3$ on MSR-VTT, $+6.5$ on VATEX), and $^{\S}$/$^{\ast}$ rows are therefore not directly
comparable in absolute terms. SCA is a \emph{single configuration} evaluated on all five
benchmarks and reported everywhere, including where it trails; $\pm$ is the standard deviation
over three seeds of that configuration. \textbf{Bold} marks the best value in each column.
``--'' not reported. AudioCaps figures for published rows are the tri-modal (T-V-A) numbers
from each paper; HyperGRAM does not report AudioCaps.}
\label{tab:main}
\small
\setlength{\tabcolsep}{5pt}
\begin{tabular}{llccccc}
\toprule
Method & Adapter & MSR-VTT & DiDeMo & ActivityNet & VATEX & AudioCaps \\
\midrule
\multicolumn{7}{l}{\emph{(a) Foundation models}} \\
ImageBind$^{\S}$        & --      & 36.8 & --   & --   & --   & \phantom{0}9.3 \\
UMT-L (25M)$^{\S}$      & --      & 40.7 & 48.6 & 41.9 & --   & -- \\
LanguageBind$^{\S}$     & --      & 44.8 & 39.9 & 41.0 & --   & 19.7 \\
mPLUG-2$^{\S}$          & --      & 47.1 & 45.7 & --   & --   & -- \\
VideoPrism-b$^{\S}$     & --      & 51.4 & --   & 49.6 & 62.5 & -- \\
VAST (27M)$^{\S}$       & full-FT & 50.7 & 49.5 & 51.4 & 82.1 & 32.1 \\
\midrule
\multicolumn{7}{l}{\emph{(b) Gramian-volume alignment}} \\
GRAM$^{\S}$                  & full-FT & 54.8 & \textbf{54.2} & \textbf{59.0} & 83.5 & 33.2 \\
GRAM, released ckpt$^{\ast}$ & full-FT & 52.5 & 50.7 & 56.3 & 90.0 & 32.2 \\
HyperGRAM$^{\S}$             & full-FT & \textbf{56.6} & 51.3 & 58.2 & 79.9 & -- \\
\midrule
\multicolumn{7}{l}{\emph{(c) Leading-eigenvalue alignment}} \\
PMRL$^{\S}$                  & full-FT & 54.5 & 50.6 & 56.0 & 80.5 & \textbf{36.1} \\
\midrule
\multicolumn{7}{l}{\emph{(d) Spherical centroid alignment (ours)}} \\
\textbf{SCA}$^{\ast}$ & LoRA
  & 54.6\tiny{$\pm$0.2}
  & 51.2\tiny{$\pm$0.3}
  & 55.9\tiny{$\pm$0.2}
  & \textbf{90.6}\tiny{$\pm$0.1}
  & 35.0\tiny{$\pm$0.4} \\
\bottomrule
\end{tabular}
\end{table*}

% READING THE TABLE -- what the numbers support, and what they do not.
%
% Against the only baseline measurable in our environment (GRAM's released checkpoint, same
% VAST foundation, same eval protocol), SCA leads on four of five:
%
%   MSR-VTT  +2.1   AudioCaps +2.8   VATEX +0.6   DiDeMo +0.5   ActivityNet -0.4
%
% MSR-VTT and AudioCaps are 5x and 7x the seed standard deviation. VATEX, DiDeMo and
% ActivityNet are within two standard deviations and should be described as level. SCA does
% this with rank-8 adapters against full fine-tuning.
%
% Against PUBLISHED numbers SCA does not lead: -2.0 on MSR-VTT and -2.3 on ActivityNet versus
% HyperGRAM, -1.1 on AudioCaps versus PMRL, and level on DiDeMo (51.2 against 51.3). Those
% gaps survive the seed error bars and must not be presented as ties. The VATEX column is not
% evidence either way -- our environment reads the released GRAM checkpoint 6.5 points above
% its published value there, so a lead over published VATEX numbers is an environment
% artefact rather than a result.
%
% The efficiency claim is separate and is measured twice: full fine-tuning is WORSE than
% rank-8 LoRA under this recipe, 52.5 (a6_fullft) and 50.9 (x3_xenc_clean) against 54.8, both
% far outside the seed spread.
